% !Mode:: "TeX:UTF-8"
\chapter*{总结\markboth{总结}{}}
\addcontentsline{toc}{chapter}{总结}

本文围绕复杂环境下UAV多模态导航的鲁棒感知问题，提出并实现了多维度可靠性感知的自适应融合框架（Idea1）。该方法通过显式可靠性估计（LiDAR SNR、RGB质量、IMU一致性）、注意力驱动的动态权重分配以及可靠性调制归一化，在特征层实现了可解释、可调试的融合路径，并与Stable-Baselines3的SAC训练流程完成端到端集成。

在工程实现方面，本文完成了所有核心模块（可靠性估计器、可靠性预测网络、动态权重分配层、可靠性调制归一化层、端到端融合模块）以及支持退化注入与难度分级的UAV多模态仿真环境，并通过集成测试验证了可运行性。当前实现中，可靠性预测网络参数量约151.8K，融合模块约719.4K；在共享特征提取器配置下，SAC策略总参数约2.35M。针对训练耗时问题，本文进一步完成了环境向量化与训练流程优化，使本地CPU下1000步训练可稳定在约32秒完成。

实验层面，本文已完成可复现对比套件，并在退化场景（degradation=0.5）下获得初步结果：中等难度下各方法均可达到较高成功率；困难难度下，无可靠性基线碰撞率上升至30\%，而本文方法保持在5\%，显示可靠性感知机制对鲁棒性具有积极作用。

后续工作将重点补齐：(1) 多随机种子统计（均值/方差与显著性检验）；(2) GPU环境下10k\textasciitilde100k步长时训练稳定性验证；(3) 更多外部代表性融合基线对比；(4) 面向真实数据与真实机平台的迁移与验证；(5) 进一步优化可靠性调制归一化的参数化形式与训练策略。
